\documentclass[10pt]{article}

\usepackage[utf8]{inputenc}

\usepackage[T1]{fontenc}

\usepackage{amsmath}

\usepackage{amsfonts}

\usepackage{amssymb}

\usepackage[version=4]{mhchem}

\usepackage{stmaryrd}

\usepackage{graphicx}

\usepackage[export]{adjustbox}

\graphicspath{ {./images/} }



\DeclareUnicodeCharacter{0131}{$\imath$}



\begin{document}

1973, т. 28, в. 4, с. $17-76 ;[11] \Pi$ о з я к Э. Г., ШІ ик и н Е. В., в кн.: Итоги науки и техники. Сер. Алгебра. Топология. Геометрия, т. 12, М., 1974, с. $171-207 ;$ [12] П о з и л к Э. Г., в кн.: с. $225-41$; [13] $\Gamma$ р и б к о в И. В., "Вестн. Моск. ун-та. Сер. Математика. Механика", 1977, Nㅜ, с, $78-83 ;[14]$ е го ж е "Успехи матем. наук", 1978, т. 33, в. 2, с. $191-92 ;[15] \mathrm{K}$ о н "Успехи матем. наук", 1978, т. 33, в. 2, с. $191-92 ;[15]$ K о н-

Ф о с с е н С. Э., там же, 1936, в. 1, c. $33-76 ;[16]$ E Ф иФ о с с е н С. Э., там же, 1936, в 1, с. $33-76 ;[16]$ E Ф и-

м о в Н. В., "Докл. АН СССР», 1961, т. $136, \mathcal{N}_{2} 6$, с. $1283-86$; [17] Е Н. В., "Докл. АН П в Н. В., П о з н я Э. Г., там же, т. 137, № 1 c. $25-27 ;[18]$ и х же, там же, № 3, с. $509-12$, [19] Е фй м о в Н. В., там же, 1962 , т. 150, № 6, с. $1206-09 ;$ [20] е г о ж е, "Матем. сб.", 1964, т. 64, Ne 2, c. $286-320$; [21] е г o же, там же, 1968 , т. 76 , № 4 , c. $499-512$; [22] III е ф е л ь С. 3., Исследован [23] е г о ж е, "Сиб. матем. Ж.», 1964, т. 5, № 6, c. 1382-96; же, "Докл. АН СССР", 1965 , т. 162 , № 2, с. $294-96$; [26] E ф им о в Н. В., "Матем. сб.", 1976, т. 100, Nㅡㄴ 3, с. 356-63; [27] е г o ж е, "Докл. АН СССР", 1953, т. 93, Nㅡ 4, с. 609-11; [28] H е i n z E., "Math. Ann.", 1955, Bd 129, H. 5, p. 451-54; [29] П o 3н я к Э. Г., "Докл. АН СССР», 1966, т. $170, N_{l} 4$, c. $786-89$; [30] е г о Ж е, "Укр. геометр. сб.", 1966, в. 3, с. 78-92; [31] Ш ик и н Е. В., "Докл. АН СССР", 1974, т. 215, № 1, c. $61-63 ;[32]$



\includegraphics[max width=\textwidth, center]{2023_07_08_caf8e092af4ded4dab85g-1}

$50-52:$ [33] B е р н е p А. Јі. “Матем. сб.", 1967, т. 74, 구 2, с. $218-40 ; 1968$, т. 75 , № 1 , c. $112-39$, т. 77, № 1 , c. $136 ;[34] \mathrm{e}$ г 0 Ж е, "Сиб. матем. Ж.», 1970 , т. 11, N1, с. $20-29$; N 4 , c. $750-$ $769 ;$ [35] е г о ж е, "Матем. заметки", 1972, т. 12, № 3, c. $281-86 ;$ [36] С о к о л о в Д. Д., "Успехи матем. наук», 1975, т. 30 , в. 1 , с. $261-62 ; 1978$, т. 33 , в. 4, с. $227-28 ; 1979$, т. 34 в. 3 , с. $213-14$; [37] А л е к с а н д р о в А. Д., Внутренняя ге ометрия выпуклых поверхностей, М.- J., 1948; [38] П о г оре п о в А. В., Виешняя геометрия выпуклых поверхностей,

М., 1969 .

H. В. Ефилов.



ОТСЧЕТА СИСТЕМА - совокупность системы координат и часов, связанных с телом, по отношению к к-рому изучается движение (или равновесие) каких-нибудь других материальных точек или тел. Любое движение является относительным, и движение тела следует рассматривать лишь по отнопению к какому-либо другому телу (телу отсчета) или системе тел. Нельзя указать, напр., как движется Луна вообще, можно Јишь определить ее движение по отношению к Земле или Солнцу и звездам и т. Д.



Математически движение тела (или материальной точки) по отношениг к выбранной О. с. описывается уравненилми, к-рые устанавливают, как изменяются с течение. времени $t$ координаты, определяющие положение тела (точки) в этой О.с. Напр., в декартовых координатах $x, y, z$ двияение точки определяется уравнениями



\[

x=f_{1}(t), y=f_{2}(t), z=f_{3}(t),

\]



ваз. уравнениями движения.



Выбор О. с. зависит от целей исследования. П ри кинематич. исследованиях все О. с. равноправны. В задачах динамики преимуцественную роль играют инерциальные системы отсчета, по отношению к которым диффференциальные уравнения имеют обычно более простой вид.



По материалам статьи Система отсчета из БСЭ-з.



ОЦЕНКА СТАТИСТИЧЕСКАЯ - ФунКция от случайных величин, применяемая для оценки неизвестных параметров төоретич. распределения вероятностей. Методы теории $\mathrm{O}$. с. служат основой современной теории ошибок; обычно в качестве неизвестных параметров выступают измеряемые физич. постоянные, а в качестве случайных величин - результаты непосредственных измерений, подверженные случайным ошибкам. Напр., если $X_{1}, X_{2}, \ldots, X_{n}$ - независимые одинаково нормально распределенные случайные величины (результаты равноточных измерений, подверженных независимым нормально распределенным случайным ошибкам), то в качестве О. с. для неизвестного среднего значения $a$ (приближенного значения измеряемой физич. постоянной) применяөтся среднее арифметическое



\[

X=\frac{1}{n}\left(X_{1}+X_{2}+\ldots+X_{n}\right) \text {. }

\]



О. с. как функция от случайных величин чаще всего задается теми шли шными формулами, выбор к-рых определяется требованиями практики. При этом различают оценки точечные и оценки интервальные.



Точечные оценки. Т о ч е ч н о й $о$ це е к о й наз. такая О.с., значение к-рой представимо геометрически в виде точки в том же иространстве, что и значения неизвестных параметров (размерность пространства равна числу оцениваемых параметров). Именно точечные О. с. и используются как приближенные значения для неизвестных физич. величин. В дальнейшем для простоты предполагается, что оценке подлөжит один единственный параметр; в этом случае точечная О.с. представляет собой функцию от результатов наблюдений, принимающую числовые значения.



Точечную О. с. наз. н е с м е щ е н н о й, если ее математич. ожидание совпадает с оцениваемой величиной, т. е. если О. с. лишена систематич. ошибки. Среднее арифметическое (1) - несмещенная О. с. для математич. ожидания одинаково распределенных случайных величин $X_{i}$ (не обязательно нормальных $)$. В то же время выборочная дисперсия



\[

\hat{s^{2}}=\frac{1}{n}\left[\left(X_{1}-\bar{X}\right)^{2}+\left(X_{2}-\bar{X}\right)^{2}+\ldots+\left(X_{n}-\bar{X}\right)^{2}\right]

\]



является смещенной $\mathrm{O}$. с. для дисперсии $\sigma^{2}=\mathrm{D} X_{i}$, так как $\hat{\hat{s}^{2}}=\left(1-\frac{1}{n}\right) \sigma^{2} ;$ в качестве несмещенной 0. с. для $\sigma^{2}$ обычно берут функцию



\[

s^{2}=\frac{n}{n-1} \hat{s^{2}}

\]



См. также Несмеценная оценка.



3 а меру точности несмещенной 0 . с. $\alpha$ для параметра $a$ чапе всего принимают дисперсию $\mathrm{D} \alpha$.



O. с. с наименьшей дисперсией наз. н а и л у ч ш е й. В приведенном примере среднее арифметическое (1) наилучшая 0 . с. Однако если распределение вероятностей случайных величин $X_{i}$ отлично от нормального, то 0. с. (1) может и не быть наилучшей. Напр., если результаты наблюдений $X_{i}$ распределены равномерно в интервале $(b, c)$, то наилучшей 0 . с. для математич. ожидания $a=(b+c) / 2$ будет полусумма крайних значений



\[

\alpha=\frac{\min X_{i}+\max X_{i}}{2} .

\]



В качестве характеристики для сравнения точности различных О. с. Ірименяют эф фе к т и в н о с т ьотношение дисперсий наилучшей оценки и данной несмещенной оценки. Напр., если результаты наблюдений $X_{i}$ распределены равномерно, то дисперсии оценок (1) и (3) выражаются формулами



\[

\mathrm{D} \bar{X}=\frac{(c-b)^{2}}{12 n}

\]



II



\[

D \alpha=\frac{(c-b)^{2}}{2(n+1)(n+2)} .

\]



Так как оценка (3) наилучшая, то эфффективвость оценки (1) в данном случае есть



\[

e_{n}(\bar{X})=6 n /(n+1)(n+2) \sim 6 / n .

\]



При большом количестве наблюдений $n$ обычно требуют. чтобы выбранная О. с. стремıлась по вероятности $\mathbf{k}$ истинному значению параметра $a$, т. е. ттобы для всякого $\varepsilon>0$



\[

\lim _{n \rightarrow \infty} P\{|\alpha-a|>\varepsilon\}=0

\]



такие О.с. наз. с о с т о я т е л ь н ы м ІІ состоятельной О. с.- люи́ая несмещенная



(пример оценка,





\end{document}